\originalchapter{电导、局域化与变程跳跃}
\label{ch:disorder}
\source{8.512第5--8讲；原PDF中部分讲号与页面讲号不一致，按主题连接。}
\originalsection{电流响应与扩散}
最小耦合$(p+eA)^2/(2m)$使电流含顺磁算符$\vct j_p$及抗磁项$-ne^2\vct A/m$。均匀场下$\vct E(\omega)=\ii\omega\vct A(\omega)$，两项相加才给规范不变电导。仅保留电流--电流相关会在正常金属错误地产生静态排磁响应。Kubo电导可写成$\sigma_{ij}(\omega)=\ii[ne^2\delta_{ij}/m+\Pi^R_{ij}(\omega)]/(\omega+\ii0^+)$，其中$\Pi^R$按电流--电流的相应符号定义。

若粒子长期运动为扩散，密度连续性与Fick定律给$\partial_t\delta n=D\nabla^2\delta n$。对化学势梯度，扩散电流和电势漂移电流在静态平衡应相抵，因而Einstein关系为$\sigma=e^2D\,\partial n/\partial\mu$。对应密度响应有扩散极点$\chi_{nn}(q,\omega)=\chi_{nn}(q,0)Dq^2/(Dq^2-\ii\omega)$。极限顺序重要：先$q\to0$且$\omega\ne0$时总粒子数无法变化，先$\omega\to0$则得到静态压缩率。

弱随机势的一次Born散射给$1/\tau=(2\pi/\hbar)D_\uparrow(E_F)\overline{|V_{kk'}|^2}$，角平均省略在记号中。输运时间乘$1-\cos\theta$，前向散射会明显区分$\tau$与$\tau_t$。在$k_F\ell\gg1$下，$D=v_F^2\tau_t/d$，恢复Drude电导。
\originalsection{边界敏感性与Thouless能量}
在边长$L$的扩散样品中，穿越时间$t_D=L^2/D$，相应能量$E_{\mathrm{Th}}=\hbar D/L^2$。单粒子平均能级间隔$\delta=1/[D(E_F)L^d]$。改变边界相位时，延展态的能量移动可达$E_{\mathrm{Th}}$；局域态的波函数在边界极小，能量移动按$\e^{-L/\xi}$下降。因此$g\sim E_{\mathrm{Th}}/\delta\sim\hbar D(E_F)DL^{d-2}$刻画电导与边界敏感性，比例常数随$e^2/h$的选法而变。

Landauer图像为$G=(e^2/h)\sum_nT_n$，自旋通道另计。理想一维通道也有两端接触电阻；四端测量与两端电导不应混称。无序相干传播需叠加振幅，不能把串联势垒的经典电阻相加当作所有尺度的精确规则。
\originalsection{局域化的尺度流}
定义$\beta(g)=\dd\log g/\dd\log L$。大$g$的经典导体给$\beta\to d-2$；局域态$g\propto\e^{-L/\xi}$给$\beta\simeq\log g$，在小$g$为负。正交对称类且无强自旋轨道与外磁场时，时间反演的一对闭合路径振幅同相，增强回返概率，降低扩散。

量子修正主要来自返回传播子$1/(Dq^2+\tau_\phi^{-1})$，所以$\delta\sigma\propto-\int_{1/L_\phi}^{1/\ell}\dd^dq/q^2$。一维修正随$L_\phi$增长，二维呈负对数，三维红外收敛。二维弱无序可在极大尺度最终局域，一维一般更快；三维可有不稳定临界点分开金属与绝缘体。以上是特定对称类的单参数尺度理论，不能推广为“所有二维体系永远绝缘”：磁场、拓扑项、自旋轨道以及相互作用都可改变流。

\begin{center}
\begin{tikzpicture}[>=Stealth]
\draw[->](0,0)--(8,0)node[right]{$\log g$};\draw[->](0,-2)--(0,1.4)node[above]{$\beta$};
\draw[thick,structurecolor](-.0,-1.8)..controls(2,-.5)and(4,1)..(7.3,1);
\draw[dashed](0,1)--(7.5,1);\node[left]at(0,1){$d-2$};
\fill(3.1,0)circle(2pt);\node[below]at(3.1,0){$g_c$};
\draw[->](1.8,-.8)--(.9,-1.3);\draw[->](4.5,.7)--(5.5,.9);
\end{tikzpicture}
\end{center}
此图仅示意三维正交类的两端极限及一个临界点，不是未经计算的数值$\beta$曲线。
\originalsection{Mott变程跳跃}
在局域绝缘体中，距$r$的两态重叠$\propto\e^{-r/\xi}$，热激活能差$\Delta$的跳跃率$\Gamma\propto\exp[-2r/\xi-\Delta/(\kb T)]$。在能窗$\Delta$与体积$r^d$中找一个态要求$\nu(E_F)\Delta r^d\sim1$。将$\Delta$消去，优化$X(r)=2r/\xi+1/[\kb T\nu r^d]$给$r_{\rm opt}^{d+1}\sim d\xi/[2\kb T\nu]$，因而
\begin{equation}
\rho(T)\sim\rho_0\exp[(T_0/T)^{1/(d+1)}],\qquad \kb T_0\sim\frac1{\nu(E_F)\xi^d}.
\end{equation}
低温最佳跳跃距离比最近邻远，故称变程。该指数要求费米能附近态密度近常数。若库仑相互作用形成软隙，优化能量--距离关系改变，常出现Efros--Shklovskii的$1/2$指数；不能用观察到$1/2$直接反推“一维Mott”。
